\section[Electronic Lab Notebook]{Electronic Lab Notebook (ELN)}\label{ssc:ELN}


\subsection{General information}

Electronic Lab Notebooks (ELNs) enable researchers to organize and store
experimental procedures, protocols, notes and data using their computer or
mobile device. ELNs can offer several advantages over the traditional paper
notebook in documenting research during the active phase of a project, including
searchability within and across notebooks, secure storage with multiple
redundancies, remote access to notebooks, and the ability to easily share
notebooks among team members and collaborators.

eLabFTW \includegraphics[height=9pt]{eLabFTW-logo3-ohneText.png} is an OpenSource, generic, browser based ELN, developed/initialized by
Nicolas Carpi (Institut Pasteuer, Paris) 2012. Data are stored in MySQL/MariaDB.
It is maintained by many developers and is used worldwide (Berkley, Indian
Institute of Technology, KIT, ...).
\begin{itemize}
  \item Access: \url{https://elabftw.iwt.zz/}
  \item Free definable status of the experiments (`finished', `running', ...)
  \item Definition of templates for a step-by-step procedure description of
        measurements
  \item Every experiment gets a unique ID and can be summarized easily as a pdf
  \item Free definable categories/tags
  \item Graphical text editor to describe your experiments and simulations
  \item File attachments of data with preview of common formats (pdf, tiff, png, …)
  \item Linking of experiments
  \item Time stamp service (RFC 3161, e.g. DFN)
  \item Data import/export (csv, zip, json, …)
  \item Access using, e.g., Python via an Application Programming Interface (API)
\end{itemize}


\subsection{IWT standard for documentation and structure}

To document scientific findings clearly and distinctly, every title of an experiment consists systematically in these six components: \\
 \\
\{Sample Number\}, \{Sample Name/Symbol\}, Material, \{Device/System\}, Method, Project Abbreviation \\
 \\
Entries with curly brackets are optional, but the three entries ``Material'', ``Method'' and ``Project'' or ``Project Abbreviation'' are abolutely imperative. You may also include more entries if required. The exactly same systematics are also applied for the tags (i.e. markers). The reason for this redundancy is the SQL database of eLabFTW which is much faster searching for tags.

You have to create a pdf for every experiment (can be a simulation as well) simply by a click (`Make a pdf’) in the edit mode (see \texttt{‘04\_PrimaryData’} in Table \ref{table:paper-directory-structure}. This kind of documentation includes a unique QR-code and is required for data which is published! The QR-code enables a direct link to the data together with unique ID.

The \textbf{maximum size} of uploaded data should not exceed
\textbf{100Mbyte per file}. Further hints and tricks may be found in the eLab-document (please add comments there if you explore new ways of usage) on the IWT file server   \\ ($\rightarrow$Austausch $\rightarrow$Forschungsdatenmanagement $\rightarrow$eLabFTW)

A template for the JSON field available in every experiment with entries for the metadata (see chapter \ref{ssc:Doku+Metadata}) is also saved on the IWT fileserver.

